\documentclass[10pt,a4paper]{amsart}
\usepackage[margin=19mm]{geometry}
\usepackage{amsmath,amssymb,mathtools}
\usepackage[scr=rsfs,cal=euler]{mathalfa}
\usepackage{xfrac}
\usepackage[unicode,hidelinks,bookmarks=false]{hyperref}
\newcommand{\ie}{i.e.\ }
\newcommand{\cf}{cf.\ }
\newcommand{\resp}{respectively\ }
\newcommand{\Subsection}{\subsection}
\providecommand{\nobreakdash}{\nobreak\hbox{-}}
\setlength{\emergencystretch}{2em}
\multlinegap0pt
\usepackage{bm}
\usepackage{bbm}
\usepackage{dsfont}
\usepackage{xspace}
\usepackage{cancel}
\usepackage{mathtools}
\usepackage{amsthm}
\makeatletter
\@ifundefined{theorem}{\newtheorem{theorem}{Theorem}}{}
\makeatother
\title[Translation Monoids and Recursive Evaluation in Finite Binar]{Translation Monoids and Recursive Evaluation in Finite Binary Algebras}
\author{Volkan Yildiz}
\date{Source: arXiv:2604.01486v1 (CC BY 4.0)}
\begin{document}
\maketitle

\noindent\emph{Source and licence.} Translation Monoids and Recursive Evaluation in Finite Binary Algebras. Version arXiv:2604.01486v1 (CC BY 4.0), distributed under the Creative Commons Attribution 4.0 International licence, as recorded in the arXiv licence metadata for this exact version and shown on the abstract page https://arxiv.org/abs/2604.01486v1. Full rights notice: licenses/arxiv-2604.01486-CC-BY-4.0.txt.\par\smallskip

\noindent\textit{Editorial scope.} Counted unit: the Theorem whose source label is \texttt{thm:context-in-translation-monoid} in the article (source0.tex lines 421--432, source label \texttt{thm:context-in-translation-monoid}), delivered together with the article's own complete original proof of it (source0.tex lines 434--483, one \texttt{proof} environment of 50 lines). No mathematics is written, repaired or re-derived: statement and proof are the article's own text line for line. Required context retained uncounted: none. Every definition, display and label of those lines is retained unchanged. Omitted from this package and not redistributed: 1--420, 433--433, 484--958. Nothing inside a retained excerpt is omitted or altered: every retained body is delivered exactly as the article's lines are written, so no omission receipt of any kind accompanies this package. Citations: the retained excerpts carry no citation command at all, so no bibliography entry of the article is transcribed and no reference is redistributed. Reference adaptation: none was needed and none was made; every reference inside the delivered unit resolves inside it, so no unresolved reference or citation is produced. Printed slips kept literal: no printed word, symbol, hypothesis or number of the counted statement or of its proof was altered. Layout and class: the article's own class is \texttt{article} with packages including geometry, amsmath, amsthm, amssymb, amsfonts, mathtools, hyperref, tikz; the wrapper is the lane's pinned \texttt{amsart} editorial wrapper at 10pt on A4, which restates the theorem-like environments the retained text needs and adds the article's own macro definitions that the retained text needs (none), with extra wrapper packages (bm, bbm, dsfont, xspace, cancel, mathtools). The delivered numbering is the unit's own. Assets: no retained span contains a figure environment, a graphics-inclusion command, a PDF-page inclusion, a table or a TikZ picture; no image file is copied into this package, the package declares no asset dependency and no image file is redistributed with it. Licence: version arXiv:2604.01486v1 of this article is distributed under the Creative Commons Attribution 4.0 International licence, as recorded in the arXiv licence metadata for this exact version and shown on the abstract page https://arxiv.org/abs/2604.01486v1. A full rights notice is delivered at licenses/arxiv-2604.01486-CC-BY-4.0.txt. Characters: every retained excerpt is pure ASCII, so the delivered engine, XeLaTeX with its default Latin Modern text font, reports no missing character.
\begin{theorem}\label{thm:context-in-translation-monoid}
Let \(A=(A,\star)\) be a finite binary algebra, let \(t\in \mathcal B_n\), and
let \(u\) be a subterm occurrence of \(t\). For every choice of values of the
variables external to \(u\), the associated context map
\[
\kappa_{t,u,\mathbf a}:A\to A
\]
belongs to the translation monoid \(T(A)\). More precisely,
\(\kappa_{t,u,\mathbf a}\) is a composition of left and right translations,
with one factor for each enclosing binary operation on the chain from \(u\) to
\(t\).
\end{theorem}

\begin{proof}
We argue by induction on the distance from \(u\) to \(t\) in the chain of
enclosing subterm occurrences.

If this distance is \(0\), then \(u=t\). In that case the context map is simply
the identity map on \(A\), since the value inserted at \(u\) is already the
value of the whole term \(t\). As \(T(A)\) is a monoid, it contains the
identity transformation, so the claim holds.

Now assume the statement holds whenever the distance is at most \(d\), and let
\(u\) be at distance \(d+1\) from \(t\). Let \(v\) be the parent subterm
occurrence of \(u\). Then exactly one of the following two cases occurs:
\[
v=s\star u
\qquad\text{or}\qquad
v=u\star s,
\]
where \(s\) is the sibling subterm occurrence of \(u\) inside \(v\).

Fix an assignment of values to all variables external to \(u\). In particular,
all variables occurring in \(s\) are fixed, so the value of \(s\) under this
assignment is a constant \(c\in A\).

Suppose first that \(v=s\star u\). Then whenever the variable input to the
context map has value \(x\in A\), the value at \(v\) is
\[
c\star x=L_c(x).
\]
From \(v\) upward to the whole term \(t\), the remaining enclosing context has
distance \(d\). By the induction hypothesis, the induced map from the value at
\(v\) to the value of \(t\) is some element \(g\in T(A)\). Hence the total
context map is
\[
\kappa_{t,u,\mathbf a}=g\circ L_c \in T(A).
\]

The case \(v=u\star s\) is analogous. The value at \(v\) is then
\[
x\star c = R_c(x),
\]
and by the induction hypothesis the remaining enclosing context contributes some
map \(g\in T(A)\). Therefore
\[
\kappa_{t,u,\mathbf a}=g\circ R_c \in T(A).
\]

In both cases the context map lies in \(T(A)\), and each step from \(u\) to
\(t\) contributes exactly one left or right translation. This completes the
induction.
\end{proof}

\end{document}
